\DocumentMetadata{
  pdfversion=2.0,
  pdfstandard=ua-2,
  lang=en-US,
  tagging=on,
  tagging-setup={math/setup=mathml-SE,tabsorder=structure}
}
\documentclass[11pt,letterpaper]{article}
\usepackage[margin=0.68in,includefoot]{geometry}
\usepackage{newcomputermodern}
\usepackage{amsmath,amscd,mathtools}
\usepackage{tikz,adjustbox}
\providecommand{\square}{\mdlgwhtsquare}
\usepackage[hidelinks]{hyperref}
\hypersetup{
  pdftitle={Algebra First-Year Exam, January 2020, Part 2},
  pdfauthor={University of Florida Department of Mathematics},
  pdfsubject={Graduate mathematics examination},
  pdfdisplaydoctitle=true,
  bookmarksopen=true
}
\setcounter{secnumdepth}{0}
\setlength{\parindent}{0pt}
\setlength{\parskip}{0.28em}
\setlength{\emergencystretch}{3em}
\setlength{\leftmargini}{2.15em}
\setlength{\leftmarginii}{2.65em}
\setlength{\leftmarginiii}{3em}
\renewcommand{\labelenumi}{\textbf{\theenumi.}}
\pagestyle{empty}
\raggedbottom
\allowdisplaybreaks
\newenvironment{examproblems}[1]{%
  \begin{enumerate}%
  \setcounter{enumi}{#1}%
  \setlength{\topsep}{0.35em}%
  \setlength{\partopsep}{0pt}%
  \setlength{\itemsep}{0.48em}%
  \setlength{\parsep}{0.18em}%
}{\end{enumerate}}
\newenvironment{examsubparts}{%
  \begin{enumerate}%
  \setlength{\topsep}{0.18em}%
  \setlength{\partopsep}{0pt}%
  \setlength{\itemsep}{0.16em}%
  \setlength{\parsep}{0.1em}%
}{\end{enumerate}}
\newenvironment{examinstructions}{%
  \par\begingroup\itshape
}{%
  \par\endgroup\vspace{0.18em}
}
\makeatletter
\renewcommand\section{\@startsection{section}{1}{0pt}%
  {-1.25ex plus -.25ex minus -.1ex}%
  {0.55ex plus .1ex}%
  {\normalfont\large\bfseries}}
\renewcommand{\@maketitle}{%
  \newpage\null\vskip -1em%
  {\footnotesize\bfseries
    \href{https://www.ufl.edu/}{University of Florida}, \href{https://math.ufl.edu/}{Department of Mathematics}\par}%
  \vskip -0.5em%
  \tagstructbegin{tag=Title}%
    {\LARGE\bfseries\href{https://gma.math.ufl.edu/exams/algebra-first-year/}{Algebra First-Year Exam}, January 2020, Part 2\par}%
  \tagstructend%
  \vskip 0.25em%
  \tagmcbegin{artifact}%
    \hrule height 1pt%
  \tagmcend%
  \vskip 1em%
}
\makeatother
\title{Algebra First-Year Exam, January 2020, Part 2}
\author{}
\date{}
\begin{document}
\maketitle
\thispagestyle{empty}
\begin{examinstructions}
Answer four problems. (If you turn in more, the first four will be graded.) Within reason, you may quote theorems as long as you state them clearly.
\end{examinstructions}
\begin{examproblems}{0}
\item
This problem has three parts.
\begin{examsubparts}
\item[\textnormal{(a)}]
Define what it means for~\(R\) to be a principal ideal domain.
\item[\textnormal{(b)}]
Give an example of an integral domain that is not a principal ideal domain.
\item[\textnormal{(c)}]
Let~\(R\) be a principal ideal domain. Suppose that, for \(i\in\mathbb{Z}\) with \(i\geq 1\), \(I_i\) is an ideal of~\(R\), and \(I_i\subseteq I_{i+1}\). Prove, from your definition of principal ideal domain, that there exists some \(n_0\in\mathbb{Z}\) with \(n_0\geq 1\) such that, for all \(n\in\mathbb{Z}\), if \(n\geq n_0\) then \(I_n=I_{n_0}\).
\end{examsubparts}
\item
Let~\(R\) be a non-zero commutative ring with identity. Prove that~\(R\) has a maximal ideal.
\item
Let~\(R\) and~\(S\) be integral domains and let \(\phi\colon R\to S\) be a unital ring homomorphism. Let \(\psi\colon R[x]\to S[x]\) be the corresponding ring homomorphism, i.e., \(\psi\) applied to a polynomial is the result of applying~\(\phi\) to each of its coefficients.
\begin{examsubparts}
\item[\textnormal{(a)}]
Let \(p(x)\in R[x]\) be non-constant and monic, and let \(q(x)=\psi(p(x))\). Assume that \(q(x)\in S[x]\) is irreducible. Prove that \(p(x)\) is irreducible in \(R[x]\).
\item[\textnormal{(b)}]
Prove that the converse need not be always true. More precisely, prove that if \(p(x)\in R[x]\) is non-constant, monic and irreducible, and we set \(q(x)=\psi(p(x))\), it need not be the case that \(q(x)\in S[x]\) is irreducible.
\end{examsubparts}
\item
Let~\(R\) be an integral domain and~\(M\) a (unital)~\(R\)-module. Suppose that~\(M\) is finitely generated, and, for each \(m\in M\), there exists some \(r\in R\) such that \(r\neq 0\) and \(rm=0\). Prove that, there exists some \(a\in R\) such that \(a\neq 0\) and, for all \(m\in M\), \(am=0\).
\item
Let \(R=\mathbb{Z}[i]\) be the ring of Gaussian integers. Prove that~\(R\) is a unique factorization domain. (You may assume standard results about different types of rings, but nothing, other than the definition and the fact that they form an integral domain, about Gaussian integers.)
\item
Let \(n\geq 1\) and \(F=\mathbb{Q}(\sqrt{1},\sqrt{2},\ldots,\sqrt{n})\). Prove that \(2^{1/3}\notin F\).
\end{examproblems}
\end{document}
